\begin{proofbox}
	\textbf{\textcolor{CProof}{思路提示}}

	通过作角平分线上的点到两边的垂线，利用角平分线性质（距离相等），利用面积比建立等式。

	\medskip
	\textbf{\textcolor{CProof}{正式推导}}
	\begin{description}[style=nextline, leftmargin=2.8em, labelsep=0.5em, itemsep=0.55em]
		\item[\textbf{步骤一（作垂线）：}]
			过点$D$作$DG\perp AB$于$G$，作$DH\perp AC$于$H$。
			\[
				DG = DH.
			\]
			\par\textbf{理由：}
			角平分线上的点到角两边距离相等。

		\item[\textbf{步骤二（面积比一）：}]
			分别以$AB$、$AC$为底，$DG$、$DH$为高表示$\triangle ABD$和$\triangle ACD$的面积。
			\[
				\begin{aligned}
					\frac{S_{\triangle ABD}}{S_{\triangle ACD}}
					&= \frac{\frac{1}{2}\cdot AB\cdot DG}{\frac{1}{2}\cdot AC\cdot DH} \\
					&= \frac{AB}{AC}.
			\end{aligned}
			\]
			\par\textbf{理由：}
			利用$DG=DH$，面积比等于两底边之比。

		\item[\textbf{步骤三（面积比二）：}]
			以$BD$、$DC$为底，过$A$作$AM\perp BC$于$M$，$AM$为公共高。
			\[
				\begin{aligned}
					\frac{S_{\triangle ABD}}{S_{\triangle ACD}}
					&= \frac{\frac{1}{2}\cdot BD\cdot AM}{\frac{1}{2}\cdot DC\cdot AM} \\
					&= \frac{BD}{DC}.
			\end{aligned}
			\]
			\par\textbf{理由：}
			两三角形同高（$A$到$BC$的垂线），面积比等于底边比。

		\item[\textbf{步骤四（等量代换）：}]
			结合上面两个面积比等式。
			\[
				\frac{AB}{AC} = \frac{BD}{DC}.
			\]
			\par\textbf{理由：}
			两个比值都等于同一个面积比，因此相等。
	\end{description}

	\medskip
	\textbf{\textcolor{CProof}{结论回扣}}

	由$\displaystyle\frac{BD}{DC}=\frac{AB}{AC}$，即$BD:DC=AB:AC$，证毕。
\end{proofbox}
